\documentclass[12pt]{article}

% Match week 2 using fonts available in the local TeX installation.
\usepackage[T1]{fontenc}
\usepackage{times}
\usepackage{microtype}
\usepackage{amsmath}
\usepackage[parfill]{parskip}
\setlength{\parindent}{0pt}
\setlength{\parskip}{2ex plus 0.2ex minus 0.2ex}
\usepackage[letterpaper,margin=1in]{geometry}
\usepackage[hidelinks]{hyperref}

\begin{document}

\begin{flushleft}
\textbf{Week 3: Stochastic Optimization} \\
Jaisel Singh (js6897) \\
September 30, 2026
\end{flushleft}

\vspace{0.1in}

What stood out to me in Hardt and Recht's optimization chapter was how much
the choice of objective matters for fitting a model. Empirical risk minimization
means minimizing average training loss, which connects directly to last week's
least-squares and likelihood objectives. Convexity helps explain why some of
these problems are tractable: for a differentiable convex objective, any
stationary point is a global minimum. Still, convergence of gradient descent
requires appropriate conditions, including the choice of step size. I found this
distinction useful because it separates properties of the problem from choices
we make when solving it.

I also found the tradeoff behind stochastic gradient descent (SGD) interesting.
For the average loss \mbox{$F(w)=n^{-1}\sum_{i=1}^{n}\ell_i(w)$}, sampling an
example uniformly gives a gradient whose expectation is $\nabla F(w)$. Each update is much cheaper
than computing the full gradient, although an individual step can increase the
training loss. This makes SGD's usefulness easier to understand: we can afford
many imperfect updates. Near an optimum, individual examples can still push
the parameters in different directions even when their gradients average to
zero. The learning-rate schedule therefore affects both how quickly we make
progress and how much we continue to fluctuate.

Gower et al.'s discussion of variance reduction gives a useful response to this
problem. For example, SVRG periodically computes a full gradient at a reference
point and uses it to correct subsequent stochastic estimates. I like the idea
of spending additional computation occasionally to make many later updates more
informative. However, Bottou et al.'s distinction between optimization and
statistical error makes me hesitate to equate a more accurate training solution
with a better model. Improving predictions also depends on the data and the
model class. The question I am left with is how, under a fixed computing budget,
we should decide whether to spend more effort reducing gradient noise or whether
ordinary SGD has already optimized the training objective well enough.

\vspace{0.05in}
{\footnotesize
\textit{Readings:} Hardt and Recht, \href{https://mlstory.org/optimization.html}{``Optimization,'' Chapter 5};
Gower et al., \href{https://arxiv.org/abs/2010.00892}{\emph{Variance-Reduced Methods for Machine Learning}};
Bottou et al., \href{https://arxiv.org/abs/1606.04838}{\emph{Optimization Methods for Large-Scale Machine Learning}}.
\par}

\end{document}

%%% Local Variables:
%%% mode: latex
%%% TeX-master: t
%%% End:
